% !PRINTED 228
\weberpage{228}
We obtain the number of dihedral groups from the fact that a dihedral group is completely determined by the cyclic group contained in it (\S.\,55). For the four-group, however, one must further observe that the same group is obtained whether one starts from $\psi$, from $\chi$, or from $\psi\chi$, so the number obtained directly must be divided by 3. Thus there are 5 groups $D_2$, 10 groups $D_3$, and 6 groups $D_5$.

Of particular importance, however, are the tetrahedral groups contained in the icosahedral group. We shall determine one of them, which we denote by $Q$.

The tetrahedral group must (by \S.\,56) contain a four-group $D_2$. We start from such a group $D_2$ and choose for it
\[
1, \psi, \chi, \psi\chi.
\]

Now further, a substitution of the third order must occur in $Q$. We may assume it in one of the two forms $\Theta^r\chi\Theta^s$ or $\Theta^r\chi\psi\Theta^s$. Since both assumptions lead to the same result, we choose as the substitution of the third order
\[
\varphi = \Theta^r\chi\Theta^s, \quad r+s\equiv\pm2 \ (\mathrm{mod}\ 5) \quad [\S.\,58, (21)].
\]

For neither of the numbers $r, s$, however, is congruence to zero ($\equiv0$) possible, since otherwise either $\chi\varphi$ or $\varphi\chi$ would be a power of $\Theta$, and hence of degree $5$, whereas no element of degree $5$ can occur in $Q$.

We further form, by \S.\,58, (23), (25):
\[
\begin{aligned}
\chi\varphi &= \Theta^r\chi\psi\Theta^{r+s}, \quad r\equiv\pm1 \\
\chi\varphi &= \Theta^{-r}\chi\Theta^{-r+s}, \quad r\equiv\pm2,
\end{aligned}
\]
whereby, for the reason cited above, $r\equiv-s$ is excluded in the first case and $r\equiv s$ in the second, and in the first case $2r+s\equiv\pm1$, in the second $-2r+s\equiv\pm2$ is required. Accordingly the four possible cases
\[
r\equiv\pm1, \quad s\equiv\pm2, \quad r\equiv\pm2, \quad s\equiv\pm1 \ (\mathrm{mod}\ 5)
\]
remain.

If one of the four substitutions of the third order so determined occurs in $Q$, then the three others occur as well. For if we set
\weq{1}{\varphi = \Theta\chi\Theta^2,}
then by \S.\,58, (23), (25), it follows that:
\weq{2}{\varphi^{-1} = \Theta^{-2}\chi\Theta^{-1}, \quad \varphi\chi = \Theta^{-1}\chi\Theta^{-2}, \quad \chi\varphi^{-1} = \Theta^2\chi\Theta.}
